\documentclass[11pt]{article}
% Preámbulo compartido por los artefactos Scrum del proyecto
% Proyecto Agile — Spanish Braille Application (EPN)
\usepackage[utf8]{inputenc}
\usepackage[T1]{fontenc}
\usepackage[spanish]{babel}
\usepackage[a4paper,margin=2.3cm]{geometry}
\usepackage{amsmath}
\usepackage{amssymb}
\usepackage{xcolor}
\usepackage{booktabs}
\usepackage{longtable}
\usepackage{array}
\usepackage{colortbl}
\usepackage{enumitem}
\usepackage{fancyhdr}
\usepackage{titlesec}
\usepackage{tcolorbox}
\usepackage{pgfplots}
\pgfplotsset{compat=1.18}
\usepackage{hyperref}

\definecolor{epnblue}{HTML}{2E4C9C}
\definecolor{epnbluelight}{HTML}{4B6FCB}
\definecolor{epnrow}{HTML}{EDF1FB}
\definecolor{epngray}{HTML}{555555}

\hypersetup{
  colorlinks=true,
  linkcolor=epnblue,
  urlcolor=epnblue,
  citecolor=epnblue,
  pdfcreator={XeLaTeX},
}

\titleformat{\section}
  {\normalfont\Large\bfseries\color{epnblue}}{\thesection}{0.6em}{}
\titleformat{\subsection}
  {\normalfont\large\bfseries\color{epnbluelight}}{\thesubsection}{0.6em}{}

\pagestyle{fancy}
\fancyhf{}
\renewcommand{\headrulewidth}{0.4pt}
\renewcommand{\footrulewidth}{0.2pt}
\fancyfoot[C]{\thepage}

\newcommand{\artefactoTitulo}[2]{%
  \begin{titlepage}
    \centering
    \vspace*{1.2cm}
    {\large Escuela Politécnica Nacional\par}
    {\normalsize Facultad de Ingeniería de Sistemas\par}
    {\normalsize Construcción y Evolución de Software --- Primer Bimestre\par}
    \vspace{1.4cm}
    \begin{tcolorbox}[colback=epnblue,colframe=epnblue,coltext=white,arc=2mm,boxrule=0pt,width=0.92\textwidth,halign=center]
      {\Huge\bfseries #1\par}
    \end{tcolorbox}
    \vspace{0.5cm}
    {\Large #2\par}
    \vspace{1.6cm}
    {\large\bfseries Proyecto Agile --- Spanish Braille Application\par}
    \vspace{0.3cm}
    {\normalsize Traducción bidireccional español \(\leftrightarrow\) Braille Unicode\par}
    \vfill
    \begin{tabular}{ll}
      \textbf{Product Owner:}  & José Castro \\
      \textbf{Scrum Master:}   & Victor Aveiga \\
      \textbf{Equipo de Desarrollo:} & Javier Angulo, Elian Caizapanta, Erick Costa, Emily Aumala \\
    \end{tabular}
    \vspace{0.8cm}
    \par\noindent\rule{0.5\textwidth}{0.4pt}\par
    \vspace{0.3cm}
    {\normalsize Documento preparado por: Erick Costa (Equipo de Desarrollo)\par}
    {\normalsize \today\par}
  \end{titlepage}
  \fancyhead[L]{Proyecto Agile --- Spanish Braille App}
  \fancyhead[R]{#1}
}


\begin{document}
\artefactoTitulo{Burndown Charts}{Seguimiento de avance por sprint}

\section{Metodología de medición}
Los gráficos de esta sección muestran la \textbf{línea ideal} (avance lineal esperado según la
capacidad y duración real de cada sprint) para el \textbf{Sprint 1} (22 jun -- 3 jul 2026, 10 días
hábiles, 12 pts) y el \textbf{Sprint 2} (6 -- 10 jul 2026, 5 días hábiles, 15 pts). Entre ambos suman los
\textbf{15 registros diarios} del Daily Scrum (pestaña 04, a cargo del Scrum Master Victor Aveiga).

Los dos puntos de la \textbf{línea real} que se pueden confirmar de forma independiente son: el inicio
(capacidad total pendiente) y el cierre de cada sprint, donde el Product Owner (José Castro) aceptó el
incremento en la Sprint Review correspondiente (0 pts pendientes al cierre). La trayectoria diaria
intermedia depende del registro día a día del tablero Kanban y del Daily Scrum, que no forma parte de
este repositorio de código; la tabla de la sección \ref{sec:registro} está lista para completarse con
esos valores reales y así reemplazar la aproximación lineal por la curva real de avance.

\section{Sprint 1 --- Transcripción Base (12 pts, 22 jun -- 3 jul)}

\begin{center}
\begin{tikzpicture}
\begin{axis}[
  width=0.95\textwidth, height=7cm,
  xlabel={Día hábil del sprint}, ylabel={Story points restantes},
  xmin=0, xmax=10, ymin=0, ymax=13,
  xtick={0,1,2,3,4,5,6,7,8,9,10}, ytick={0,3,6,9,12},
  grid=both, grid style={line width=.1pt, draw=gray!25},
  major grid style={line width=.2pt,draw=gray!45},
  legend pos=north east, legend style={font=\small},
  axis lines=left,
]
\addplot[thick, dashed, color=epngray, mark=none] coordinates {(0,12) (2,9.6) (4,7.2) (6,4.8) (8,2.4) (10,0)};
\addlegendentry{Ideal}
\addplot[thick, color=epnblue, mark=*, mark size=1.8pt, mark options={fill=epnblue}] coordinates {(0,12) (10,0)};
\addlegendentry{Real (confirmado: inicio/cierre)}
\end{axis}
\end{tikzpicture}
\end{center}

\section{Sprint 2 --- Funcionalidades Avanzadas (15 pts, 6 -- 10 jul)}

\begin{center}
\begin{tikzpicture}
\begin{axis}[
  width=0.95\textwidth, height=7cm,
  xlabel={Día hábil del sprint}, ylabel={Story points restantes},
  xmin=0, xmax=5, ymin=0, ymax=16,
  xtick={0,1,2,3,4,5}, ytick={0,3,6,9,12,15},
  grid=both, grid style={line width=.1pt, draw=gray!25},
  major grid style={line width=.2pt,draw=gray!45},
  legend pos=north east, legend style={font=\small},
  axis lines=left,
]
\addplot[thick, dashed, color=epngray, mark=none] coordinates {(0,15) (1,12) (2,9) (3,6) (4,3) (5,0)};
\addlegendentry{Ideal}
\addplot[thick, color=epnblue, mark=*, mark size=1.8pt, mark options={fill=epnblue}] coordinates {(0,15) (5,0)};
\addlegendentry{Real (confirmado: inicio/cierre)}
\end{axis}
\end{tikzpicture}
\end{center}

\noindent\textbf{Sprint 2 es un time-box mucho más corto} (5 días hábiles frente a los 10 del Sprint 1)
con mayor capacidad (15 vs 12 pts), lo que implica una velocidad diaria objetivo casi el triple
(3 pts/día vs 1.2 pts/día). Esto es relevante para la exposición del Scrum Master sobre la salud del
sprint: el equipo debió sostener un ritmo de entrega notablemente más alto en la segunda mitad del
proyecto.

\section{Registro diario para completar (Scrum Master)}
\label{sec:registro}
Tabla lista para transcribir los valores reales día a día desde el tablero Kanban / Daily Scrum
(pestaña 04). El valor \textit{Ideal} ya está calculado; la columna \textit{Real} debe completarse antes
de la presentación final.

\begin{center}
\renewcommand{\arraystretch}{1.2}
\begin{tabular}{|c|c|c|c|}
\hline
\rowcolor{epnblue}
\textcolor{white}{\textbf{Sprint}} & \textcolor{white}{\textbf{Día}} & \textcolor{white}{\textbf{Ideal (pts)}} & \textcolor{white}{\textbf{Real (pts)}} \\
\hline
1 & 0 & 12.0 & 12 \\
\hline
1 & 2 & 9.6 & \rule{1.2cm}{0.4pt} \\
\hline
1 & 4 & 7.2 & \rule{1.2cm}{0.4pt} \\
\hline
1 & 6 & 4.8 & \rule{1.2cm}{0.4pt} \\
\hline
1 & 8 & 2.4 & \rule{1.2cm}{0.4pt} \\
\hline
1 & 10 & 0.0 & 0 \\
\hline
2 & 0 & 15.0 & 15 \\
\hline
2 & 1 & 12.0 & \rule{1.2cm}{0.4pt} \\
\hline
2 & 2 & 9.0 & \rule{1.2cm}{0.4pt} \\
\hline
2 & 3 & 6.0 & \rule{1.2cm}{0.4pt} \\
\hline
2 & 4 & 3.0 & \rule{1.2cm}{0.4pt} \\
\hline
2 & 5 & 0.0 & 0 \\
\hline
\end{tabular}
\end{center}

\section{Qué dicen estos datos sobre la salud del sprint}
\begin{itemize}[leftmargin=1.4cm]
  \item Ambos sprints cerraron con el Incremento aceptado en la Sprint Review, es decir, 0 pts
        pendientes al final del time-box en los dos casos.
  \item El Sprint Backlog (documento anterior) registra 32 h técnicas en Sprint 1 y 23 h en Sprint 2,
        repartidas entre los 4 integrantes del Equipo de Desarrollo sin jerarquías internas.
  \item El único punto de atención documentado por el equipo es la exportación PNG/PDF de 06H6
        (Sprint 2): reportada como completada por el equipo, pero pendiente de confirmar en el código
        de este repositorio (ver nota en el Sprint Backlog).
\end{itemize}

\end{document}
